\documentclass[12pt]{report}                            % font size and type of document
\usepackage[a4paper, margin=25mm]{geometry}             % paper size and margins
\usepackage[onehalfspacing]{setspace}                   % line spacing
\usepackage{lipsum}                                     % create dummy text in order to check formating

\usepackage[ngerman]{babel}                             % load multilingual support for german
\usepackage[babel,german=guillemets]{csquotes}          % ensure multilingual support for biblatex quoting

\usepackage{iftex}
\ifPDFTeX
   \usepackage[utf8]{inputenc}
   \usepackage[T1]{fontenc}                             % font encoding for pdflatex compiler
   \usepackage{lmodern}
\else
   \ifXeTeX
     \usepackage{fontspec}                              % font encoding for xelatex compiler
     \usepackage{unicode-math}                          % math encoding and math fonts support
     \setmathfont{Cambria Math}
   \else 
     \usepackage{luatextra}                             % font encoding for lulatex compiler
   \fi
   \setmainfont{Trebuchet MS}
   \defaultfontfeatures{Ligatures=TeX}
\fi

\newcommand{\titlename}{Latex Bedienungsanleitung}  % change title of document here
\newcommand{\version}{0.1}                                 % change version number here
\newcommand{\authorname}{Daniel Klein-Günnewick}            % change author name here
\newcommand{\state}{Verbindlich, in Bearbeitung}            % change document state here
\newcommand{\lasteditor}{David Efing}                       % change last editor name here
\newcommand{\lastchanged}{15.02.2024 09:05}                 % change last editor name here
\newcommand{\company}{Lebbing automation \& drives GmbH}    % change company name her

\usepackage[nostruts]{titlesec}                             % enable editing of chapter, section, ... titles
\titleformat{\chapter}[block]
  {\normalfont\LARGE\bfseries}{\thechapter}{10pt}{}         % customize chapter titles
\titlespacing*{\chapter}{0pt}{-30pt}{5pt}
\titleformat{\section}
  {\normalfont\Large\bfseries}{\thesection}{10pt}{}         % customize section titles
\titlespacing*{\section}{0pt}{0pt}{5pt}
\titleformat{\subsection}
  {\normalfont\large\bfseries}{\thesubsection}{10pt}{}      % customize subsection title
\titlespacing*{\subsection}{0pt}{0pt}{5pt}
\titleformat{\subsubsection}
  {\normalfont\normalsize\bfseries}{\thesubsubsection}{10pt}{}
\titlespacing*{\subsubsection}{0pt}{0pt}{5pt}               % customize subsubsection title

\usepackage[titles]{tocloft}                            % extension for toc, lof, lot and other lists
\renewcommand{\cftchapleader}{\cftdotfill{\cftdotsep}}  % dotted chapter leaders
\renewcommand{\cftchapfont}{\normalfont}                % toc entries in normal font style
\renewcommand{\cftchappagefont}{\normalfont}            % toc page numbers in normal font style
\cftsetindents{figure}{0pt}{40pt}                       % control space to left margin and the width of the numbers in lof
\cftsetindents{table}{0pt}{40pt}                        % control space to left margin and the width of the numbers in lot
\setlength{\cftchapnumwidth}{20pt}                      % adjust the width of the chapter numbers
\setlength{\cftsecnumwidth}{30pt}                       % adjust the width of the section numbers
\setlength{\cftsubsecnumwidth}{40pt}                    % adjust the width of the subsection numbers
%\setlength{\cftfignumwidth}{40pt}                       % adjust the width of the figure numbers
%\setlength{\cfttabnumwidth}{40pt}                       % adjust the width of the table numbers
\usepackage{tocbibind}                                  % adds all lists to table of contents

\usepackage{fancyhdr}                                   % customize header and footer
\makeatletter
\fancypagestyle{plain}{
  \fancyhf{}                                              % delete all header and footer content from page style
  \fancyhead[L]{\includegraphics[width=4.0cm]{LebbingLogo.png}}   % left header content
  \fancyhead[R]{\titlename~ V\version \par 
  Status: \state 
  \par \nouppercase{\leftmark}}                                   % right header content
  \setlength{\headheight}{43pt}                           % height of header for more space
  \fancyfoot[L]{\company}                                 % company name in left footer
  \fancyfoot[R]{\thepage}                                 % page number in right footer
}
\pagestyle{fancy}                                       % define new page style called fancy
\fancyhf{}                                              % delete all header and footer content from page style
\fancyhead[L]{\includegraphics[width=4.0cm]{LebbingLogo.png}}   % left header content
\fancyhead[R]{ \titlename~ V\version \par 
Status: \state 
\par \nouppercase{\leftmark}}                           % right header content
\setlength{\headheight}{43pt}                           % height of header for more space
\fancyfoot[L]{\company}                                 % company name in left footer
\fancyfoot[R]{\thepage}                                 % page number in right footer
\renewcommand{\headrulewidth}{0pt}                      % delete horizontal line in header
\renewcommand{\chaptermark}[1]{
\markboth{\thechapter{} #1}{}
}
\makeatother

\title{\titlename}                                      % for pdf document information
\author{\authorname}
\date{\today}
\newcounter{savepage}                                   % save page number inside document

\usepackage{xurl}                                       % hyphening links
\usepackage{hyperref}
\hypersetup{
    colorlinks=true,
    linkcolor=black,
    citecolor=black,
    filecolor=magenta,      
    urlcolor=cyan,
    pdftitle={\titlename},
    pdfauthor={\authorname}
    }
\urlstyle{same}                                         % allows links within pdf document and to external websites

\usepackage{graphicx}                                   % required for inserting images
\graphicspath{{images/}}
\usepackage{wrapfig}                                    % wrap text around figures
\usepackage{caption}                                    % more customization of captions
%\captionsetup{figurename=Abb., tablename=Tab.}
\DeclareCaptionType{mycapequ}[][Formelverzeichnis]      % adds list of equations
\DeclareCaptionLabelFormat{nan}{Gleichung #2}           % caption compatibility with equations 
\captionsetup[mycapequ]{labelformat=nan}

\usepackage{xcolor}                                     % load xcolor package for customization
\usepackage{listings}                                   % supports listings environments
\definecolor{codegreen}{rgb}{0,0.6,0}                   % define and set custom colors
\definecolor{codegray}{rgb}{0.5,0.5,0.5}
\definecolor{codepurple}{rgb}{0.58,0,0.82}
\definecolor{backcolour}{rgb}{0.95,0.95,0.92}
\lstdefinestyle{mystyle}{                               % define listings style
    language={C++},
    morecomment=[l]{//},
    morekeywords={String},
    backgroundcolor=\color{backcolour},   
    commentstyle=\color{codegreen},
    keywordstyle=\color{magenta},
    numberstyle=\tiny\color{codegray},
    stringstyle=\color{codepurple},
    basicstyle=\ttfamily\footnotesize,
    breakatwhitespace=false,         
    breaklines=true,                 
    captionpos=b,                    
    keepspaces=true,                 
    numbers=left,                    
    numbersep=5pt,                  
    showspaces=false,                
    showstringspaces=false,
    showtabs=false,                  
    tabsize=2
}
\renewcommand{\lstlistlistingname}{Quellcodeverzeichnis}
\renewcommand{\lstlistingname}{Quellcode}               % change displayed name of list of listings
\lstset{style=mystyle}                                  % set listings style to my defined style

\usepackage[acronym,toc,translate=babel,
            nonumberlist, nopostdot]{glossaries}        % adds support for glossary and acronyms and abbreviations
\renewcommand{\glsnamefont}[1]{\textbf{#1}}             % glossary headings in bold style
\setlength\glsdescwidth{1.0\hsize}                      % adjust the width of the acronym and glossary lists
\makenoidxglossaries
% add new acronyms here
\newacronym{tia}{TIA}{Totally Integrated Automation}                                    % input acronyms.tex and glossary.tex files
% add new glossary entries here
\newglossaryentry{latex}{name=LaTeX,description={ist ein plattformunabhängiges und freies Softwarepaket, das die Benutzung des Textsatzsystems TeX mit Hilfe von Makros vereinfacht. LaTeX ist eine Auszeichnungssprache und ein Dateiformat für digitale Texte, insbesondere für Texte, die mathematische Formeln enthalten. LaTeX wurde Anfang der 1980er Jahre von Leslie Lamport entwickelt. Der Name bedeutet so viel wie Lamport TeX. Die Entwicklung wurde seit den 1990er Jahren von einer Anzahl Entwicklern weitergeführt. Heute ist LaTeX die populärste Methode, TeX zu verwenden. Das Softwarepaket liegt mit Stand vom Dezember 2023 in der Version LaTeX release 2023-11-01 vor.}}

\usepackage[ngerman, noabbrev]{cleveref}                % improves referencing of tables, figures, equations, etc.
%\crefformat{equation}{Gl.~(#2#1#3)}                    % customise reference to Gl. (default "Gleichung")
\crefname{listing}{Quellcode}{Quellcodes}               % change reference names
\Crefname{listing}{Quellcode}{Quellcodes}

\usepackage[final]{pdfpages}                            % adds support for pdf file integration
\usepackage{tikz}                                       % create graphics in latex environment
\usetikzlibrary{positioning, calc}                      % enable relative positioning of nodes in tikz figure
\usepackage{pgfplots}                                   % create diagram, charts etc. with tikz library
\pgfplotsset{compat=1.18}                               % set used pgfplot package version for compatibility
\usepackage{amsmath}                                    % enhanced mathematical expressions
\usepackage{siunitx}                                    % includes si units support
\newcommand{\nmsbracket}[3]{\ensuremath{{}^{\phantom{#2}#1}\mkern-#3mu\mbox{}_{\phantom{#1}#2}\mathopen{[}}}
%\usepackage{mathtools}                                  % more math symbols
%\usepackage{tabulary}                                   % advanced table options
%\usepackage{makecell}
\usepackage{longtable}                                  % for large tables that should continue on the next pages
\newcommand{\colwidth}[1]{\dimexpr#1\linewidth-2\tabcolsep-1.5\arrayrulewidth}
%\usepackage{array}
%\setlength{\tabcolsep}{10pt}                            % control the space from text to cell borders
%\renewcommand{\arraystretch}{1.35}                      % set the cell height
\usepackage{tabto}                                      % tabulator support
\usepackage{enumitem}                                   % for customizing itemize/enumerate
\newcommand{\multilineitemize}[1]{                      % custom command for controlling itemize environments inside tables
  \begin{minipage}[t]{\linewidth}
    \begin{itemize}[nosep,left=0pt,label=--,after=\strut]
      #1
    \end{itemize}
  \end{minipage}
}

\usepackage[                                            % enable bibliography
  backend=biber, 
  natbib=true,
  hyperref=true,
  style=ieee
]{biblatex}
%\DeclareLanguageMapping{ngerman}{german-apa}           % only necessary for biblatex-apa style
%\DeclareFieldInputHandler{extradate}{\def\NewValue{}}  % remove suffix letters if author and year
\addbibresource{references.bib}                         % are identical in bibliography

\usepackage{afterpage}                                  % commands get expanded after the current page  
\newcommand\blankpage{                                  % is output using \afterpage{\command}
    \newpage                                            % -> useful for \clearpage
    \null                                               % define command to insert blank page
    \thispagestyle{empty}
    \newpage}
